\section{Summary}

In this chapter we introduced

\begin{itemize}
	\item limits;
	\item derivatives;
	\item higher derivatives;
	\item partial derivatives;
	\item total derivatives;
	\item the differential operator $\nabla$;
	\item the gradient;
	\item divergence;
	\item curl;
	\item the Laplacian.
\end{itemize}

These operators form the mathematical foundation of nearly every field of
theoretical physics. In the next chapter we shall apply them to classical
mechanics and derive Newton's laws of motion in a mathematically rigorous way.


%%%%%%%%%%%%%%%%%%%%%%%%%%%%%%%%%%%%%%%%%%%%%%%%%%%%%%%%%%%%%%%%%%%%%%%%%%%%%%%%

\clearpage
\chapterformulaheading
\begin{formulasheetbox}
\begin{formulatable}
Derivative & $f'(x)=\displaystyle\lim_{h\to0}\frac{f(x+h)-f(x)}{h}$\\
Gradient & $\nabla f=
\mathbf e_x\partial_x f+\mathbf e_y\partial_y f+\mathbf e_z\partial_z f$\\
Divergence & $\nabla\cdot\mathbf A=
\partial_xA_x+\partial_yA_y+\partial_zA_z$\\
Curl & $\nabla\times\mathbf A$\\
Laplacian & $\nabla^2 f=
\partial_x^2f+\partial_y^2f+\partial_z^2f$\\
Total derivative & $\displaystyle\frac{df}{dt}
=\partial_t f+\dot x\,\partial_xf+\dot y\,\partial_yf+\dot z\,\partial_zf$\\
\end{formulatable}
\end{formulasheetbox}

\chapternotationheading
\begin{notationbox}
\begin{notationtable}
$d/dx$ & Ordinary derivative operator & inverse unit of $x$\\
$\partial/\partial x$ & Partial derivative operator & inverse unit of $x$\\
$\nabla$ & Vector differential operator & m$^{-1}$ in space\\
$\nabla f$ & Gradient of scalar field & unit of $f$ per m\\
$\nabla\cdot\mathbf A$ & Divergence of vector field & unit of $\mathbf A$ per m\\
$\nabla\times\mathbf A$ & Curl of vector field & unit of $\mathbf A$ per m\\
$\nabla^2$ & Laplacian & m$^{-2}$ in space\\
\end{notationtable}
\end{notationbox}

\begin{physical}
The gradient points toward the steepest local increase of a scalar field.
Divergence measures local source or sink behaviour. Curl measures local
circulation. The Laplacian compares a value with its spatial neighbourhood.
These interpretations reappear in electrostatics, fluid mechanics, diffusion,
and the Schrödinger equation.
\end{physical}

\begin{warningbox}
A partial derivative changes one independent variable while holding the others
fixed. A total derivative follows the complete dependence of a quantity,
including indirect dependence through variables such as $x(t)$, $y(t)$,
and $z(t)$.
\end{warningbox}

\begin{chapterhistorical}[Newton and Leibniz]
Isaac Newton and Gottfried Wilhelm Leibniz developed calculus independently in
the seventeenth century. Newton emphasized changing physical quantities, while
Leibniz introduced notation that remains central to modern analysis.
\end{chapterhistorical}

\chapteroutcomesheading
\begin{outcomesbox}
The reader should now be able to:
\begin{itemize}
\item define a derivative using a limit;
\item distinguish ordinary, partial, and total derivatives;
\item compute gradients, divergence, curl, and Laplacians;
\item give geometric and physical interpretations of these operators;
\item recognize the same operators in classical and quantum field equations.
\end{itemize}
\end{outcomesbox}

\chapterexercisesheading
\subsection*{Basic}
\begin{chapterexercise}[Derivative from first principles]\difficulty{2}
Use the limit definition to derive $d(x^2)/dx=2x$.
\end{chapterexercise}

\begin{chapterexercise}[Gradient]\difficulty{1}
Calculate $\nabla f$ for $f(x,y,z)=x^2y+yz^2$.
\end{chapterexercise}

\subsection*{Intermediate}
\begin{chapterexercise}[Divergence and curl]\difficulty{2}
For $\mathbf A=(xy,yz,zx)$, calculate
$\nabla\cdot\mathbf A$ and $\nabla\times\mathbf A$.
\end{chapterexercise}

\subsection*{Advanced}
\begin{chapterexercise}[Laplacian and waves]\difficulty{3}
Calculate $\nabla^2 e^{i\mathbf k\cdot\mathbf r}$ and interpret the result.
\end{chapterexercise}

\chapterlookaheadheading
\begin{lookingaheadbox}
The next chapter applies derivatives to momentum and acceleration. Kinematics
becomes dynamics when forces determine how motion changes:
\[
\mathbf a
\longrightarrow
m\mathbf a
\longrightarrow
\mathbf F.
\]
\end{lookingaheadbox}
